\begin{theorem} \label{thm:t1}
    Sean $A\subseteq\Rn{n}$ un conjuto abierto y conexo y  $f:A\subseteq\Rn{n}\to\R$ un campo de clase $\mathcal{C}^1(\interior{A})$. Si $\boldsymbol{\sigma}:[a,b]\to A\subseteq\Rn{n}$ es una trayectoria continua,  entonces
    \[
       \int_{\boldsymbol{\sigma}}\grad f\cdot d\mathbf{s}=f(\boldsymbol{\sigma}(b))-f(\boldsymbol{\sigma}(a)).
    \]    

    \textit{Demostración.}
    Al aplicar la regla de la cadena a la función compuesta
    \begin{equation*}
        F:t \rightarrow f(\boldsymbol{\sigma}(t))
    \end{equation*}
    se obtiene
    \begin{equation*}
        F'(t)=(f\circ\boldsymbol{\sigma})'(t)=\grad{f(\boldsymbol{\sigma})}\cdot\boldsymbol{\sigma}'(t)
    \end{equation*}
    La función F es una función con valores reales de la variable $t$ y, por lo tanto, por el teorema fundamental
    del cálculo de una variable:
    \begin{equation*}
        \int_{a}^{b}F'(t)dt=F(b)-F(a)=f(\boldsymbol{\sigma}(b))-f(\boldsymbol{\sigma}(a)).
    \end{equation*}
    Por consiguiente,
    \begin{align*}
        \int_{\boldsymbol{\sigma}}\grad f\cdot d\mathbf{s}&=\int_{a}^{b}\grad{f(\boldsymbol{\sigma}(t))}\cdot\boldsymbol{\sigma}'(t)dt\\
        &=\int_{a}^{b}F'(t)dt=F(b)-F(a)\\
        &=f(\boldsymbol{\sigma}(b))-f(\boldsymbol{\sigma}(a)).
    \end{align*}
\end{theorem}



\begin{corollary}\label{tmh:c1}
    Sean $A\subseteq\Rn{n}$ un conjuto abierto y conexo y  $f:A\subseteq\Rn{n}\to\R$ un campo de clase $\mathcal{C}^1(\interior{A})$.  Si  $\boldsymbol{\sigma}_1:[a_1,b_1]\to A_1\subseteq\Rn{n}$ y $\boldsymbol{\sigma}_2:[a_2,b_2]\to A_2\subseteq\Rn{n}$ son dos trayectorias continuas  tales que $\boldsymbol{\sigma}_1(a_1)=\boldsymbol{\sigma}_2(a_2)\;\land\;\boldsymbol{\sigma}_1(b_1)=\boldsymbol{\sigma}_2(b_2)$, entonces
    \[
    \int_{\boldsymbol{\sigma}_1}\grad f\cdot d\mathbf{s}=\int_{\boldsymbol{\sigma}_2}\grad f\cdot d\mathbf{s}.
    \]
    
 Es decir,  $\int_{\boldsymbol{\sigma}}\grad f\cdot d\mathbf{s}$  s\'olo depende de los puntos inicial y final de la trayectoria.

\textit{Demostración.} 
Usando el Teorema \ref{thm:t1}, se tiene que
\begin{align*}
    \int_{\boldsymbol{\sigma}_1}\grad f\cdot d\mathbf{s}&=f(\boldsymbol{\sigma}_1(b_1))-f(\boldsymbol{\sigma}_1(a_1))\\
    \int_{\boldsymbol{\sigma}_2}\grad f\cdot d\mathbf{s}&=f(\boldsymbol{\sigma}_2(b_2))-f(\boldsymbol{\sigma}_2(a_2)).
\end{align*}
Así, con $\boldsymbol{\sigma}_1(a_1)=\boldsymbol{\sigma}_2(a_2)\;\land\;\boldsymbol{\sigma}_1(b_1)=\boldsymbol{\sigma}_2(b_2)$,
\begin{equation*}
    f(\boldsymbol{\sigma}_1(a_1))=f(\boldsymbol{\sigma}_2(a_2))\;\land\;f(\boldsymbol{\sigma}_1(b_1))=f(\boldsymbol{\sigma}_2(b_2))
\end{equation*}
Luego,
\begin{equation*}
    f(\boldsymbol{\sigma}_1(b_1))-f(\boldsymbol{\sigma}_1(a_1))=f(\boldsymbol{\sigma}_2(b_2))-f(\boldsymbol{\sigma}_2(a_2))
\end{equation*}
Por lo tanto, se demuestra que:
\begin{equation*}
    \int_{\boldsymbol{\sigma}_1}\grad f\cdot d\mathbf{s}=\int_{\boldsymbol{\sigma}_2}\grad f\cdot d\mathbf{s}.
\end{equation*}
\end{corollary}


\begin{corollary}\label{tmh:c2}
    Sean $A\subseteq\Rn{n}$ abierto y conexo y $f:A\subseteq\Rn{n}\to\R$ de clase $\mathcal{C}^1$. Si $\mathbf{F}:A\subseteq\Rn{n}\to\Rn{n}$ es un campo vectorial para el cual $\exists f:A\subseteq\Rn{n}\to\R$ clase $\mathcal{C}^1$ tal que 
    $$\grad f(\mathbf{x})=\mathbf{F}(\mathbf{x})\quad\forall\:\mathbf{x}\in A,$$
    entonces para todo par de curvas $C_1$ y $C_2$ en $A$, con los mismos extremos y misma orientaci\'on,
    \[
    \int_{C_1}\mathbf{F}\cdot d\mathbf{s}=\int_{C_2}\mathbf{F}\cdot d\mathbf{s}. 
    \]
    Dada esta igualdad, se dice que la integral de línea de un campo conservativo tiene \textit{independencia del camino}.

    \textit{Demostración.} 
    Sean $C_1$ y $C_2$ parametrizadas por las trayectorias $\boldsymbol{\sigma}_1:[a_1,b_1]\to C_1$ y $\boldsymbol{\sigma}_2:[a_2,b_2]\to C_2$, con $\grad f(\mathbf{x})=\mathbf{F}(\mathbf{x})$:
    \begin{align*}
        \int_{C_1}\mathbf{F}\cdot d\mathbf{s}&=\int_{\boldsymbol{\sigma}_1}\grad f\cdot d\mathbf{s}\\
        \int_{C_2}\mathbf{F}\cdot d\mathbf{s}&=\int_{\boldsymbol{\sigma}_2}\grad f\cdot d\mathbf{s}
    \end{align*}
    Además, como $C_1$ y $C_2$ comparten extremos, se da que $\boldsymbol{\sigma}_1(a_1)=\boldsymbol{\sigma}_2(a_2)\;\land\;\boldsymbol{\sigma}_1(b_1)=\boldsymbol{\sigma}_2(b_2)$.
    Por lo tanto, según el Corolario \ref{tmh:c1}:
    \begin{align*}
        \int_{\boldsymbol{\sigma}_1}\grad f\cdot d\mathbf{s}&=\int_{\boldsymbol{\sigma}_2}\grad f\cdot d\mathbf{s}\\
        \therefore \int_{C_1}\mathbf{F}\cdot d\mathbf{s}&=\int_{C_2}\mathbf{F}\cdot d\mathbf{s}
    \end{align*}
\end{corollary}

\begin{corollary}\label{tmh:c3}
    Sean $A\subseteq\Rn{n}$ abierto y conexo y $f:A\subseteq\Rn{n}\to\R$ de clase $\mathcal{C}^1$. Si $\mathbf{F}=\grad f$, entonces
    $$\oint_C\mathbf{F}\cdot d\mathbf{s}=0,$$
    para toda curva $C$ simple cerrada contenida en $A$.

    \textit{Demostración.} 
    Sea $C$ parametrizada por la trayectoria continua $\boldsymbol{\sigma}:[a,b]\to C$, y $\mathbf{F}=\grad f$:
    \begin{equation*}
        \oint_C\mathbf{F}\cdot d\mathbf{s}=\int_{\sigma}\grad f\cdot d\mathbf{s}
    \end{equation*}
    Cumpliéndose las hipótesis del Teorema \ref{thm:t1},
    \begin{equation*}
        \int_{\sigma}\grad f\cdot d\mathbf{s}=f(\boldsymbol{\sigma}(b))-f(\boldsymbol{\sigma}(a))
    \end{equation*}
    Ahora, como $C$ es cerrada, $\boldsymbol{\sigma}(a)=\boldsymbol{\sigma}(a)$, entonces $f(\boldsymbol{\sigma}(b))=f(\boldsymbol{\sigma}(a))$, y,
    \begin{align*}
        \int_{\sigma}\grad f\cdot d\mathbf{s}&=0\\
         \therefore \oint_C\mathbf{F}\cdot d\mathbf{s}&=0
    \end{align*}

\end{corollary}

\textit{Notación:} $\int_{\mathbf{x}_0}^{\mathbf{x}}\mathbf{F}\cdot d\mathbf{s}$ indica la integral de camino de el campo $\mathbf{F}$ sobre la recta que une los puntos $\mathbf{x}_0$
    y $\mathbf{x}$, orientada desde el primero hacia el segundo.

\begin{theorem}\label{tmh:t2}
   Sea  $\mathbf{F}:A\subseteq\Rn{n} \to \Rn{n}$ continuo con  $A$ abierto y conexo. Si  $\int_C \mathbf{F}\cdot d\mathbf{s}$ es independiente del camino en $A$  entonces el campo
    $$f:A\subseteq\Rn{n}\to\R\: | \:f(\mathbf{x})=\int_{\mathbf{x}_0}^{\mathbf{x}}\mathbf{F}\cdot d\mathbf{s},  \mbox{ con }  \mathbf{x}_0\in A, $$
    cumple que: 
    \begin{enumerate}
        \item $f$ es $\mathcal{C}^1$.
        \item $\grad f(\mathbf{x})=\mathbf{F}(\mathbf{x})\quad\forall\:\mathbf{x}\in A.$
    \end{enumerate}
\end{theorem}

\begin{theorem}
    Sean $A\subseteq\Rn{n}$ abierto y conexo, $\mathbf{F}:A\subseteq\Rn{n}\to\R$ continuo. Las siguientes proposiciones sobre $\mathbf{F}$ son equivalentes:
    \begin{enumerate}
        \item $\exists f:A\subseteq\Rn{n}\to\R$ de clase $\mathcal{C}^1$ tal que $\mathbf{F}(\mathbf{x})=\grad f(\mathbf{x}),\;\forall\mathbf{x}\in A.$ 
        \item La integral de $\mathbf{F}$ es independiente del camino en $A$.
        \item $\oint_C\mathbf{F}\cdot d\mathbf{s}=0$ para toda curva cerrada simple $C$ en $A$.
    \end{enumerate}

    \textit{Demostración.}
    Se deben demostrar las implicaciones de a pares. 
    
    \begin{equation}
        \mathbb{F}(\mathbf{x})=\grad f(\mathbf{x}) \then \int_{C_1}\mathbf{F}\cdot d\mathbf{s}=\int_{C_2}\mathbf{F}\cdot d\mathbf{s}.
        \label{eq:CampoCons1}
    \end{equation}
    Con $C_1$ y $C_2$ de iguales extremos. En tanto $\mathbb{F}(\mathbf{x})=\grad f(\mathbf{x})$, y $\boldsymbol{\sigma}_1:[a_1,b_1]\to C_1$ una parametrización continua de $C_1$, según el Teorema \ref{thm:t1},
    \begin{equation*}
        \int_{C_1}\mathbf{F}\cdot d\mathbf{s}=\int_{\boldsymbol{\sigma}_1}\grad f(\mathbf{x})\cdot d\mathbf{s}=f(\boldsymbol{\sigma}_1(b_1))-f(\boldsymbol{\sigma}_1(a_1))
    \end{equation*}
    Esto se cumple análogamente para $C_2$ con parametrización $\boldsymbol{\sigma}_2:[a_2,b_2]\to C_2$. Como $C_1$ y $C_2$ comparten extremos
    \begin{align*}
        \boldsymbol{\sigma}_1(a_1)=\boldsymbol{\sigma}_2(a_2) \quad &\land \quad \boldsymbol{\sigma}_1(b_1)=\boldsymbol{\sigma}_2(b_2)\\
        f(\boldsymbol{\sigma}_1(a_1))=f(\boldsymbol{\sigma}_2(a_2)) &\quad \land \quad f(\boldsymbol{\sigma}_1(b_1))=f(\boldsymbol{\sigma}_2(b_2))\\
        \then f(\boldsymbol{\sigma}_1(b_1))-f(\boldsymbol{\sigma}_1(a_1))&=f(\boldsymbol{\sigma}_2(b_2))-f(\boldsymbol{\sigma}_2(a_2))\\
        \therefore \int_{C_1}\mathbf{F}\cdot d\mathbf{s}&=\int_{C_2}\mathbf{F}\cdot d\mathbf{s}
    \end{align*}

    \begin{equation}
        \int_{C_1}\mathbf{F}\cdot d\mathbf{s}=\int_{C_2}\mathbf{F}\cdot d\mathbf{s} \then \mathbb{F}(\mathbf{x})=\grad f(\mathbf{x}).
        \label{eq:CampoCons2}
    \end{equation}
    Esta proposición es verdadera por el Teorema \ref{tmh:t2}.

    A partir de las ecuaciones (\ref{eq:CampoCons1}) y (\ref{eq:CampoCons2}) se demostró que
    \begin{equation}
        \mathbb{F}(\mathbf{x})=\grad f(\mathbf{x}) \Leftrightarrow \int_{C_1}\mathbf{F}\cdot d\mathbf{s}=\int_{C_2}\mathbf{F}\cdot d\mathbf{s}.
        \label{eq:CampoCons1y2}
    \end{equation}

    \begin{equation}
        \mathbb{F}(\mathbf{x})=\grad f(\mathbf{x}) \then \oint_C\mathbf{F}\cdot d\mathbf{s}=0
        \label{eq:CampoCons3}
    \end{equation}
    Esta proposición es verdadera por el Corolario \ref{tmh:c3}.

    \begin{equation}
        \oint_C\mathbf{F}\cdot d\mathbf{s}=0 \then \int_{C_1}\mathbf{F}\cdot d\mathbf{s}=\int_{C_2}\mathbf{F}\cdot d\mathbf{s}
        \label{eq:CampoCons4}
    \end{equation}
    Para la demostración pensemos en dos curvas con iguales extremos, $C_1$ y $C_2$.
    \begin{figure}[H]
        \begin{center}
            \begin{tikzpicture}[scale=2, >=stealth]

                % Puntos A y B
                \coordinate (A) at (0,0);
                \coordinate (B) at (3,0);

                % Primer camino: curva hacia arriba (gamma_1)
                \draw[thick, blue, ->, decoration={markings, mark=at position 0.5 with {\arrow{>}}},
                    postaction={decorate}] 
                    (A) .. controls (1,1.5) and (2,1.5) .. (B) node[midway, above] {$C_1$};

                % Segundo camino: curva hacia abajo (gamma_2)
                \draw[thick, red, ->, decoration={markings, mark=at position 0.5 with {\arrow{>}}},
                    postaction={decorate}] 
                    (B) .. controls (1,-1.5) and (2,-1.5) .. (A) node[midway, below] {$C_2$};

                % Nombres de los puntos
                \filldraw (A) circle(0.03) node[below left] {$A$};
                \filldraw (B) circle(0.03) node[below right] {$B$};

            \end{tikzpicture}
            \caption{Curvas $C_1$ y $C_2$.}
        \end{center}
    \end{figure}
    Luego, se puede pensar la curva $C=C_1\cup\overline{C_2}$ con $\overline{C_2}$ la curva $C_2$ en sentido opuesto. Así:
    \begin{align*}
        \oint_C\mathbf{F}\cdot d\mathbf{s}&=\int_{C_1}\mathbf{F}\cdot d\mathbf{s}+\int_{\overline{C_2}}\mathbf{F}\cdot d\mathbf{s}\\
        0&=\int_{C_1}\mathbf{F}\cdot d\mathbf{s}-\int_{C_2}\mathbf{F}\cdot d\mathbf{s}\\
        \therefore \int_{C_1}\mathbf{F}\cdot d\mathbf{s}&=\int_{C_2}\mathbf{F}\cdot d\mathbf{s}
    \end{align*}
   
    A partir de haber demostrado (\ref{eq:CampoCons3}), (\ref{eq:CampoCons4}) y (\ref{eq:CampoCons1y2}) se construye lógicamente las equivalencias de este teorema.
\end{theorem}

